\centerline{\bf \Huge Основная теорема арифметики}


\begin{flushright}
        \scriptsize{Episode 1: Angel Attack}
\end{flushright}

\textbf{Основная теорема арифметики.} Любое натуральное число $n \geq 2$ представимо единственным образом с точностью до перестановки сомножителей в виде произведения $n=p_1^{\alpha_1} p_2^{\alpha_2} \ldots p_k^{\alpha_k}$, где $p_1, p_2, \ldots, p_k$ различные простые числа.

\noindent
\textbf{Важная лемма.} Если $(a,c)=1$ и $c\ |\ ab$, то $c\ | \ b$. 

\problem{1}
Докажите иррациональность числа (а) $\sqrt{2}$; (б) $\sqrt[3]{\frac{16}{33}}$.

\problem{2}
Даны натуральные числа $a$ и $b$, причём $a<1000$. Докажите, что если $a^{21}$ делится на $b^{10}$, то $a^2$ делится на $b$.

\problem{3}
Найдите все такие натуральные $n$, для которых число $n^2-1$ является степенью простого числа.

\problem{4}
Можно ли расставить 2025 натуральных чисел по кругу так, чтобы отношение любых двух соседних было простым числом?

\problem{5 }
В клетки таблицы размером $45 \times 45$ расставили все натуральные числа от 1 до 2025. Вычислили произведения чисел в каждой строке таблицы и получили набор из 45 чисел. Затем вычислили произведения чисел в каждом столбце таблицы и также получили набор из 45 чисел. Могли ли полученные наборы оказаться одинаковыми?

\problem{6}
Натуральные числа $x, y$ таковы, что оба числа $x^3+y$ и $y^3+x$ делятся на $x^2+y^2$.

(a) Докажите, что числа $x$ и $y$ взаимно просты. \textit{Подсказка:} пусть НОД$(x,y) = d$

(б) Найдите все такие числа $x$ и $y$.
%https://problems.ru/view_problem_details_new.php?id=98440

\problem{7}
Найдите все натуральные числа $x$ и $y$, для которых $x^y=y^x$.
%https://problems.ru/view_problem_details_new.php?id=77873
